\documentclass[11pt,a4paper]{article}
\usepackage[utf8]{inputenc}
\usepackage[T1]{fontenc}
\usepackage[english]{babel}
\usepackage{amsmath,amssymb,mathtools}
\usepackage{geometry}
\usepackage{enumitem}
\usepackage{hyperref}
\geometry{margin=2.4cm}
\hypersetup{colorlinks=true,linkcolor=blue,urlcolor=blue}

\title{Mathematical Formulation: Multimodal Patent Classification via Semantic Distillation, Concept Lattice, and Organizational Inference}
\author{}
\date{\today}

\begin{document}
\maketitle

\begin{abstract}
This document formulates a research architecture for classifying patents into Monitoring (M), Prediction (P), Recommendation (R), and Intervention (I). The system receives the patent title, abstract, and figure descriptions generated by a Vision-Language Model (VLM). A Large Language Model (LLM) proposes structured semantic cues; a student model learns from existing labels and teacher annotations. A Concept Lattice is built from extracted attributes and used at inference time to compute category scores. The resulting cues are passed to a separate stage for inferring company capabilities and organizational structure.
\end{abstract}

\section{Definitions and Notation}
Let:
\begin{itemize}
\item $\mathcal{P}=\{p_1,\ldots,p_N\}$ is the set of training patents;
\item $x_p=(t_p,a_p,v_p)$ is the patent input, where $t_p$ is the title, $a_p$ is the abstract, and $v_p$ is the textual description of figures generated by the VLM;
\item $y_p\in\mathcal{Y}=\{M,P,R,I\}$ is the existing patent label;
\item $c(p)$ is the company associated with the patent, retained as metadata for the organizational stage and not used as input to the M/P/R/I classifier;
\item $\mathcal{E}_{c}$ is the external evidence available about company $c$;
\item $\mathcal{A}$ is the universe of candidate semantic attributes discovered automatically from the corpus.
\end{itemize}
No manually defined ontology is assumed. Candidate attributes may include terms, phrases, semantic roles, actions, outcomes, and relations extracted by the teacher.

\section{Multimodal Representation}
A text encoder $f_{\theta}$ represents the title and abstract, while $g_{\phi}$ represents the textual description produced by the VLM:
\begin{align}
h_p^T &= f_{\theta}([t_p;a_p]),\\
h_p^V &= g_{\phi}(v_p),\\
z_p &= F_{\psi}(h_p^T,h_p^V).
\end{align}
$F_{\psi}$ is a multimodal fusion module, for example concatenation followed by projection or cross-attention.

\section{Semantic Teacher and Cue Extraction}
The LLM teacher $T_{\omega}$ receives $x_p$ and generates structured cues:
\begin{equation}
Q_p=T_{\omega}(x_p)=\{q_{p1},\ldots,q_{pm_p}\}.
\end{equation}
Each cue is represented as:
\begin{equation}
q_{pi}=(s_{pi},r_{pi},\ell_{pi},e_{pi},u_{pi}),
\end{equation}
where $s_{pi}$ is the term or phrase, $r_{pi}$ the semantic role, $\ell_{pi}$ the source location (text span or figure, when available), $e_{pi}$ the relations to other cues, and $u_{pi}$ the uncertainty. Cues are semantic hypotheses grounded in the input, not automatic facts about the company.

For training, $Q_p$ can be generated offline and audited on a sample. At inference time, the same annotation schema is used, either by the teacher or by the distilled student.

\section{Teacher--Student Distillation}
The student $S_{\eta}$ produces:
\begin{equation}
S_{\eta}(x_p)=(\hat{\boldsymbol{\pi}}_p,\hat{Q}_p,\hat{R}_p),
\end{equation}
where $\hat{\boldsymbol{\pi}}_p$ is the class distribution, $\hat{Q}_p$ the extracted cues, and $\hat{R}_p$ the relations between cues. The class distribution can be computed as:
\begin{equation}
\hat{\boldsymbol{\pi}}_p=\operatorname{softmax}(W_yz_p+b_y).
\end{equation}
If the teacher produces a calibrated distribution $\boldsymbol{\pi}^{T}_p$, distillation is defined as:
\begin{equation}
\mathcal{L}_{KD}^{cls}=\tau^2\sum_{p\in\mathcal{P}}
D_{KL}(\boldsymbol{\pi}^{T}_{p,\tau}\Vert
\hat{\boldsymbol{\pi}}_{p,\tau}),
\end{equation}
where $\tau$ is the distillation temperature. If the teacher does not produce calibrated probabilities, this term may be omitted.

A baseline student loss is:
\begin{align}
\mathcal{L}_{student}={}&
\lambda_{cls}\mathcal{L}_{CE}
+\lambda_{KD}\mathcal{L}_{KD}^{cls}\\
&+\lambda_{span}\mathcal{L}_{span}
+\lambda_{rel}\mathcal{L}_{rel}.
\end{align}
$\mathcal{L}_{CE}$ is cross-entropy for the existing labels; $\mathcal{L}_{span}$ trains span extraction, and $\mathcal{L}_{rel}$ trains relation extraction. Teacher annotations may be noisy and should be audited.

\section{Building the Concept Lattice}

\subsection{Formal Context}
We define the formal context as:
\begin{equation}
\mathbb{K}=(\mathcal{P},\mathcal{A},\mathcal{R}),
\end{equation}
where $(p,a)\in\mathcal{R}$ indicates that attribute $a$ is associated with patent $p$. In classical binary FCA:
\begin{equation}
B_{pa}=\mathbf{1}[ (p,a)\in\mathcal{R}].
\end{equation}
For $X\subseteq\mathcal{P}$ and $Y\subseteq\mathcal{A}$, the derivation operators are:
\begin{align}
X'&=\{a\in\mathcal{A}:\forall p\in X,\ (p,a)\in\mathcal{R}\},\\
Y'&=\{p\in\mathcal{P}:\forall a\in Y,\ (p,a)\in\mathcal{R}\}.
\end{align}
A formal concept is a pair $(X,Y)$ satisfying $X'=Y$ and $Y'=X$. Concepts ordered by inclusion of their extents form a lattice:
\begin{equation}
(X_1,Y_1)\preceq(X_2,Y_2)
\iff X_1\subseteq X_2
\iff Y_2\subseteq Y_1.
\end{equation}
The lattice is derived from attributes extracted from the corpus. M/P/R/I labels are used to learn concept relevance, not to impose a manually specified ontology. A weighted variant with values in $[0,1]$ is possible, but it is an extension of classical binary FCA and must be defined separately.

\subsection{Associating Concepts with Categories}
Let $C$ be the set of concepts in the lattice. For concept $c\in C$ and category $k\in\mathcal{Y}$, define:
\begin{equation}
\rho(c,k)=
\frac{\sum_{p\in\mathcal{P}}\mathbf{1}[y_p=k]\mu(p,c)}
{\sum_{p\in\mathcal{P}}\mu(p,c)+\epsilon},
\end{equation}
where $\mu(p,c)$ measures the association between patent $p$ and concept $c$, and $\epsilon>0$. A concept may have high values for multiple categories; this overlap is retained explicitly.

\section{Scoring a New Patent}
For a new patent $p^*$, the student produces $\hat{Q}_{p^*}$. Its match to concept $c$ is $\mu(p^*,c)$, and $w(c)$ is a learned or experimentally validated concept weight. The category score is:
\begin{equation}
s_k(p^*)=
\frac{\sum_{c\in C}\mu(p^*,c)w(c)\rho(c,k)}
{\sum_{c\in C}\mu(p^*,c)w(c)+\epsilon},
\quad k\in\mathcal{Y}.
\end{equation}
The calibrated scores are:
\begin{equation}
\pi_k(p^*)=\frac{\exp(s_k(p^*)/T)}
{\sum_{j\in\mathcal{Y}}\exp(s_j(p^*)/T)},
\end{equation}
where $T>0$ is the calibration temperature. For mutually exclusive classification:
\begin{equation}
\hat{y}_{p^*}=\arg\max_{k\in\mathcal{Y}}\pi_k(p^*).
\end{equation}
If the labeling scheme allows overlapping labels, multi-label classification should be used instead of selecting a single maximum.

\section{Explainable Evidence and Provenance}
The raw contribution of concept $c$ to the score of class $k$ is:
\begin{equation}
\Delta_{c,k}(p^*)=\mu(p^*,c)w(c)\rho(c,k).
\end{equation}
This provides a structural explanation of the score, but is not automatically a causal explanation of the model. It can be validated using ablations or occlusion tests.

Each cue passed to the organizational stage retains:
\begin{equation}
d_i=(\text{concept},\text{surs\u{a}},\text{loca\u{t}ie},
\text{scor},\text{incertitudine},\text{tipul dovezii}).
\end{equation}
The evidence type distinguishes an explicit statement in the patent, a semantic interpretation, and a hypothesis about company capability.

\section{Organizational Inference}
The company $c$ associated with the patent is known, and $\mathcal{E}_c$ may contain external evidence about it. These data are not used for M/P/R/I classification. For company $c$, we construct:
\begin{equation}
G_c=(V_c,E_c,X_c),
\end{equation}
where nodes represent candidate capabilities or organizational entities, edges represent relations between them, and $X_c$ contains their attributes. The model estimates:
\begin{equation}
P(G_c\mid D_c,\mathcal{E}_c),
\end{equation}
where $D_c$ represents evidence extracted from the company's patents. Without direct evidence, the graph is a probabilistic inference, not a confirmed organizational chart; every node and edge should retain its source and uncertainty.

\section{Joint Objective and Training}
One possible multi-objective formulation is:
\begin{align}
\mathcal{L}={}&
\lambda_1\mathcal{L}_{cls}
+\lambda_2\mathcal{L}_{KD}
+\lambda_3\mathcal{L}_{span}
+\lambda_4\mathcal{L}_{rel}\\
&+\lambda_5\mathcal{L}_{lattice}
+\lambda_6\mathcal{L}_{org}.
\end{align}
$\mathcal{L}_{lattice}$ must be defined concretely, for example to reward concept--category associations that generalize on validation data. $\mathcal{L}_{org}$ should be used only if reference graphs or verifiable organizational evidence are available. These losses should not be assigned values without data and an operational definition.

Classical FCA produces a discrete structure; backpropagation through lattice construction requires an explicit differentiable relaxation or an alternating procedure. A robust initial plan is to train the student, build the lattice from stabilized attributes, and then tune concept scoring.

\section{Evaluation}
Evaluate the following separately:
\begin{enumerate}
\item M/P/R/I classification: macro-F1, balanced accuracy, and calibration;
\item cue extraction: span-level precision, recall, and F1;
\item semantic relations: precision/recall/F1 on validated relations;
\item lattice: stability, size, and benefit compared with a classifier without a lattice;
\item organizational graph: node and edge quality where reference data exist.
\end{enumerate}
Data splits must prevent leakage between patent families; testing on companies absent from training is also recommended. Teacher annotations should be audited to estimate noise.

\section{Limitations}
\begin{itemize}
\item The operational definitions of M/P/R/I come from the existing labeling scheme; they are not invented by the model.
\item VLM-generated descriptions may contain errors, and their provenance must be retained.
\item The presence of a concept in a patent does not prove that the company has an internal team implementing that capability.
\item Company identity is used in the organizational stage, not as a classification shortcut.
\item This formulation does not establish novelty or patentability; a prior-art search and analysis of relevant claims are required.
\end{itemize}
\end{document}